\subsection{Translational Requirements Derived from the Evidence}

The purpose of the present section is not to introduce additional biological
claims, but to apply the evidence-to-engineering procedure defined in
Section~2 to the findings established in Sections~3--5. The central
translational question is whether the experimental conditions under which TMR
has demonstrated measurable effects can be represented as system-level
requirements without converting unresolved biological or technical assumptions
into apparent capabilities.
The seven general requirement categories defined in Section~2 are consolidated
here into five intervention-specific requirements because measurement quality,
intervention control, operational reliability, and empirical validation span
multiple TMR-specific functions rather than forming independent implementation
requirements.

Five requirements follow consistently from the preceding evidence.

First, the intervention requires \textbf{traceable memory preparation}. TMR
depends on an association established before sleep, but the literature does not
support a universal pre-sleep performance threshold below which cueing becomes
invalid. Prior memory strength and the structure of cue--memory associations
influence subsequent TMR effects in a condition-dependent manner
\citep{creery2015priorlearning,cairney2016benefits}. A first-generation system
should therefore preserve the identity of cue--memory mappings and record
pre-sleep learning performance, while treating any numerical eligibility
threshold as a parameter requiring prospective validation.

Second, the intervention requires \textbf{physiological eligibility}. Cue
presentation should be restricted to sleep conditions supported by the TMR
literature rather than governed solely by elapsed time. Meta-analytic evidence
supports significant average effects during NREM stage~2 and slow-wave sleep,
although neither state is universally optimal across all paradigms
\citep{hu2020promoting}. An autonomous system must therefore estimate whether
the current physiological interval is eligible for stimulation and must be
capable of withholding intervention when that judgment is unsupported.

Third, the system requires \textbf{conservative disturbance control}. TMR
depends on preserving the sleep physiology within which cueing occurs.
Unsupervised and automated home studies indicate that stimulus intensity and
sleep disturbance can materially influence behavioral outcome
\citep{goldi2019home,whitmore2022improving}. Cue delivery should therefore be
governed by signal quality, physiological stability, recent stimulation
history, and evidence of arousal or state transition. The literature does not
establish universal thresholds for these variables; those thresholds remain
implementation-specific validation questions.

Fourth, the intervention requires \textbf{temporal validity}. Physiological
information used to authorize cueing must remain relevant when the cue is
actually delivered. The system must therefore measure the delay between signal
acquisition, state estimation, decision logic, scheduling, and physical cue
onset. No single maximum acceptable latency can be derived from the present
evidence because the relevant timescale depends on whether control is based on
broad sleep stage, transient events, or oscillatory phase. End-to-end latency
should consequently be treated as a measurable system property whose acceptable
range must be established relative to the physiological target.

Fifth, experimental validation requires \textbf{intervention traceability}.
Delivered and withheld cues should be reconstructable in relation to the
physiological evidence available at the time of the decision. At minimum, the
validation system should retain synchronized information concerning cue
identity, physiological state estimate, signal quality or confidence,
decision time, actual cue onset, and subsequent evidence of sleep disruption.
Behavioral outcomes can then be interpreted in relation to how the system
actually operated rather than according to its intended protocol alone.

These requirements do not imply that every variable must ultimately appear in
a deployed product. Some, particularly detailed intervention logging and
cued--uncued behavioral comparisons, are requirements of scientific validation
rather than permanent user-facing functionality. The purpose of the
first-generation system is to make the intervention experimentally
interpretable before optimization for usability or scale.

Table~\ref{tab:translation_requirements} summarizes the translation from
scientific evidence to system requirement while retaining the principal
uncertainty associated with each step.

\begin{table*}[t]

\centering

\caption{
Evidence-derived requirements for a first-generation closed-loop TMR system.
The final column identifies assumptions that remain to be established
prospectively rather than treated as settled design parameters.
}

\label{tab:translation_requirements}

\footnotesize
\setlength{\tabcolsep}{4.5pt}
\renewcommand{\arraystretch}{1.22}

\begin{tabularx}{\textwidth}{
@{}
>{\raggedright\arraybackslash}p{2.7cm}
>{\raggedright\arraybackslash}X
>{\raggedright\arraybackslash}X
>{\raggedright\arraybackslash}X
@{}
}

\toprule

\textbf{Evidence basis}
&
\textbf{Translational implication}
&
\textbf{System-level requirement}
&
\textbf{Unresolved validation question}
\\

\midrule

\textbf{Cue--memory association}
&
TMR operates on information associated with sensory cues before sleep.
&
Preserve cue--memory identity and record pre-sleep learning performance.
&
What encoding or retrieval criterion, if any, should determine eligibility for
overnight cueing?
\\

\addlinespace[1.2mm]

\textbf{State-dependent efficacy}
&
Cueing effects depend on the physiological context in which stimulation occurs.
&
Estimate intervention eligibility causally and permit abstention during
unsupported or unstable states.
&
What state definition and confidence criterion provide adequate intervention
reliability?
\\

\addlinespace[1.2mm]

\textbf{Sensitivity to disturbance}
&
Excessive or poorly controlled stimulation may disrupt sleep and reduce or
reverse behavioral benefit.
&
Control cue exposure and reassess physiological stability following
stimulation.
&
Which post-cue signals and thresholds best identify biologically consequential
disturbance?
\\

\addlinespace[1.2mm]

\textbf{Temporal dependence}
&
The physiological evidence authorizing stimulation may become outdated before
cue onset.
&
Measure and synchronize the complete sensing-to-stimulation pathway.
&
What end-to-end delay is acceptable for the selected level of physiological
targeting?
\\

\addlinespace[1.2mm]

\textbf{Condition-dependent behavioral effect}
&
System operation must be interpretable in relation to actual cueing conditions.
&
Maintain synchronized intervention records during experimental validation.
&
Which operational variables best predict successful versus unsuccessful TMR
trials?
\\

\bottomrule

\end{tabularx}

\end{table*}


\subsection{Comparison of First-Generation Implementation Pathways}

The requirements above constrain the implementation space but do not uniquely
identify a device or architecture. Selection therefore requires comparison of
candidate pathways according to two distinct questions: how strongly each
pathway is connected to the existing evidence, and how many additional
assumptions must be validated before an unsuccessful or successful result can
be interpreted.

\noindent\textbf{Fixed-schedule versus physiology-guided cueing.}

The relevant distinction is not between all ``open-loop'' and ``closed-loop''
TMR experiments. Many foundational studies achieved physiological control
through researcher-supervised PSG rather than through an automated feedback
algorithm. The translational problem arises when an autonomous system delivers
cues according to elapsed time or a predetermined schedule without verifying
that the sleeper remains in an appropriate state.

For an unsupervised system, physiology-guided control provides substantially
greater interpretability. Current-state information can determine whether a cue
is authorized, delayed, or withheld, while post-cue information can determine
whether stimulation should continue. This introduces additional sensing and
control complexity, but that complexity serves a scientific function: it
preserves the relationship between the intervention and the physiological
conditions under which TMR efficacy has been demonstrated. Closed-loop control
is therefore selected here as a requirement of \textit{autonomous
state-dependent implementation}, not as a claim that TMR itself is invalid
without an automated closed loop.


\noindent\textbf{Wearable EEG versus peripheral physiological inference.}

The sensing comparison requires particular caution because the available
evidence supports both approaches, but at different levels of physiological
observability.

Wearable EEG retains access to scalp electrophysiology used to characterize
sleep state and selected events relevant to TMR. Reduced-montage systems have
demonstrated substantial sleep-staging capability, although performance and
signal usability vary across devices, populations, and recording conditions
\citep{arnal2020dreem,ravindran2025dreem}. More importantly for the present
translation problem, wearable EEG has already been incorporated into automated
home TMR: Salfi et al.\ used real-time EEG-derived slow-wave detection to guide
vocabulary cueing and reported a behavioral advantage for cued material
\citep{salfi2025promoting}.

Peripheral sensing should not be characterized as experimentally unvalidated.
PPG and related signals can provide substantial information about sleep state
\citep{korkalainen2020ppg}, and SleepStim demonstrated that smartwatch-derived
cardiac and inertial information can govern automated TMR cue delivery at home
\citep{whitmore2022improving}. Its behavioral results also demonstrated,
however, the importance of stimulus intensity and intervention conditions.
The current evidence therefore establishes peripheral closed-loop feasibility,
but does not establish that peripheral-state estimates provide intervention
control equivalent to electrophysiological monitoring across TMR protocols.

No direct head-to-head experiment currently demonstrates that an EEG-guided TMR
controller produces greater behavioral efficacy than an otherwise equivalent
peripheral controller. The choice between these pathways must therefore remain a
translational inference rather than an experimentally established ranking.
EEG offers closer continuity with the electrophysiological measurements used in
much of the controlled TMR literature, whereas peripheral sensing reduces
measurement burden at the cost of an additional inference between measured
autonomic physiology and the cortical state relevant to the intervention.


\noindent\textbf{Stage-aware versus event- or phase-aware targeting.}

The literature also supports different levels of temporal specificity.
Meta-analytic evidence demonstrates TMR benefits when cueing is restricted to
broad NREM conditions, particularly N2 and slow-wave sleep
\citep{hu2020promoting}. At finer temporal scales, spindle-related findings and
closed-loop experiments indicate that the momentary oscillatory context of cue
presentation can influence subsequent memory
\citep{antony2018spindle,shimizu2018closedloop,ngo2022shaping,
nicolas2025phase}.

Phase-specific control therefore has direct experimental support and should not
be treated merely as a mechanistically plausible future extension. However,
the available evidence remains narrower across memory domains, protocols, and
independent replications than the evidence supporting TMR within broader NREM
conditions. Moreover, existing phase-specific studies do not establish a
universal incremental advantage over well-controlled stage-aware TMR. For
example, Nicolas et al.\ found superior performance for UP- relative to
DOWN-state cueing, but the UP-reactivated condition did not significantly
outperform the non-reactivated condition \citep{nicolas2025phase}.

Greater temporal specificity also introduces additional dependencies on
real-time event detection, phase estimation, causal processing,
synchronization, and physical cue timing. For the initial reference pathway,
the relevant distinction is therefore not between a scientifically valid and
an invalid strategy, but between two evidence-supported strategies with
different levels of implementation dependency.

The first-generation pathway is consequently defined here as stage-aware and
non-phase-locked. This choice should be interpreted as a sequencing decision,
not as evidence that phase-aware control is ineffective or unnecessary.
Fine-scale slow-wave or spindle information may be recorded during
first-generation validation, while the incremental value of phase-aware control
should subsequently be tested against the characterized reference
configuration.


\noindent\textbf{Existing-platform integration versus proprietary hardware.}

The decision to begin with existing wearable EEG hardware follows a different
type of reasoning. It is not a biological conclusion but an engineering
sequencing judgment. Developing new sensing hardware simultaneously with a new
closed-loop intervention would introduce additional uncertainty concerning
electrode design, acquisition quality, synchronization, comfort, durability,
and hardware reliability. Failure of the integrated experiment could then be
difficult to localize.

Using existing hardware narrows this first validation problem only if the
platform exposes the functions required by the intervention. A candidate device
must provide sufficiently usable physiological signals in real time, reliable
timestamps or synchronization, access compatible with causal processing, and a
means of coordinating cue delivery with the recorded physiology. If no
available platform satisfies these conditions, existing-platform integration
cannot be assumed to be preferable simply because the hardware already exists.

The initial platform strategy should therefore be stated conditionally:
development should preferentially integrate with a compatible and sufficiently
validated wearable EEG platform rather than begin with proprietary hardware.
The specific device remains an engineering-selection question outside the
evidentiary conclusion of this paper.


\begin{table*}[t]

\centering

\caption{
Comparison of candidate implementation choices for a first-generation
closed-loop TMR reference system. The table separates demonstrated evidence
from uncertainties that would accompany each pathway.
}

\label{tab:implementation_tradeoffs}

\footnotesize
\setlength{\tabcolsep}{4.5pt}
\renewcommand{\arraystretch}{1.22}

\begin{tabularx}{\textwidth}{
@{}
>{\raggedright\arraybackslash}p{2.5cm}
>{\raggedright\arraybackslash}X
>{\raggedright\arraybackslash}X
>{\raggedright\arraybackslash}X
@{}
}

\toprule

\textbf{Decision axis}
&
\textbf{Candidate options}
&
\textbf{Evidence-supported interpretation}
&
\textbf{First-generation implication}
\\

\midrule

\textbf{Cue control}
&
Fixed-schedule autonomous cueing versus physiology-guided closed-loop control.
&
TMR efficacy is state-dependent; autonomous fixed scheduling cannot verify the
physiology present at cue onset.
&
Use physiology-guided control with explicit ability to withhold stimulation.
\\

\addlinespace[1.2mm]

\textbf{Sensing}
&
Peripheral physiological inference versus wearable EEG.
&
Both have participated in automated home TMR. Peripheral sensing infers neural
state indirectly; EEG preserves scalp electrophysiological observability
\citep{whitmore2022improving,salfi2025promoting}.
&
Use wearable EEG as the initial reference pathway while retaining peripheral
control as an experimentally viable alternative.
\\

\addlinespace[1.2mm]

\textbf{Temporal targeting}
&
Stage-aware versus event- or phase-aware cueing.
&
NREM stage-dependent efficacy is established more broadly than any incremental
benefit of precise phase targeting
\citep{hu2020promoting,antony2018spindle}.
&
Begin with stage-aware, non-phase-locked control; test finer temporal targeting
subsequently.
\\

\addlinespace[1.2mm]

\textbf{Platform}
&
Compatible existing wearable EEG versus newly developed hardware.
&
No biological evidence requires proprietary hardware; new hardware would add
an independent acquisition and validation problem.
&
Prefer existing-platform integration when real-time signal access,
synchronization, and intervention control are technically available.
\\

\bottomrule

\end{tabularx}

\vspace{1mm}

\begin{minipage}{\textwidth}
\footnotesize
\textit{Note.}
The table does not rank technologies by universal quality. It identifies the
pathway that introduces the fewest additional assumptions into the initial
test of the closed-loop TMR intervention.
\end{minipage}

\end{table*}


\subsection{Evidence-Based Selection of the First-Generation Reference Pathway}

The preceding comparison does not demonstrate that one architecture is
universally superior. No study identified in the present synthesis directly
compares complete EEG-guided and peripheral-guided TMR systems under identical
participants, learning tasks, cueing policies, and outcome measures. The
first-generation decision must therefore be understood as an
\textit{evidence-based translational judgment}: the preferred pathway is the one
that allows the central intervention hypothesis to be tested while introducing
the fewest additional unvalidated dependencies.

On this basis, the most defensible first-generation reference pathway is an
\textbf{EEG-guided, stage-aware, real-time closed-loop TMR software layer
intended for integration with compatible existing wearable EEG hardware}.
This selection follows from the convergence of several observations rather than
from any single study.

Controlled TMR evidence establishes that cueing efficacy depends on
physiological state and produces a modest, condition-dependent behavioral effect
\citep{hu2020promoting}. Wearable EEG has demonstrated that substantial
sleep-related electrophysiological information can be retained in reduced,
portable systems, while also revealing practical limitations related to signal
quality and population generalization
\citep{arnal2020dreem,ravindran2025dreem}. Most directly, EEG-guided automated
TMR has produced a measurable vocabulary benefit under home conditions
\citep{salfi2025promoting}. Together, these findings provide an empirical chain
from laboratory physiology to wearable measurement to automated intervention.

Peripheral sensing remains scientifically credible. Sleep-stage inference from
cardiac and multimodal peripheral signals is well supported, and SleepStim
provides direct evidence that peripheral physiology can control automated home
TMR \citep{korkalainen2020ppg,whitmore2022improving}. The reason for not
selecting it as the initial reference is therefore not absence of feasibility.
Instead, a peripheral-first system would require the first validation program
to establish both the TMR intervention and the adequacy of an indirect
physiological control signal at the same time. EEG retains greater
electrophysiological observability and allows failures in cueing, sleep-state
detection, and post-stimulation response to be investigated more directly.

Stage-aware rather than phase-locked control follows the same sequencing
principle. The behavioral evidence for NREM-targeted TMR is broader than the
evidence that precise oscillatory-phase targeting yields additional behavioral
benefit \citep{hu2020promoting}. Beginning with non-phase-locked control reduces
the number of timing assumptions that must be validated simultaneously.
Fine-scale slow-wave or spindle information remains scientifically valuable and
may be recorded during first-generation experiments, but it should become a
mandatory control variable only if prospective comparison demonstrates an
advantage over the simpler reference policy.

Existing-platform integration similarly reduces the number of independent
technical problems introduced into the first experiment. This choice remains
conditional rather than absolute: an existing EEG device is suitable only if
its signal access, timing behavior, data quality, and integration interfaces can
satisfy the requirements established above. The paper therefore selects a
\textit{platform strategy}, not a specific commercial device.

The resulting decision can be expressed as a sequence of controlled
uncertainty:

\begin{quote}
    validated TMR evidence $\rightarrow$ electrophysiological state observation $\rightarrow$ causal eligibility decision $\rightarrow$ controlled cue delivery $\rightarrow$ post-cue observation $\rightarrow$ behavioral validation.
\end{quote}

The value of the selected pathway is that each link can be observed and tested
without simultaneously requiring a new sensing modality, a phase-sensitive
controller, and a newly developed hardware platform to succeed.

This selection is deliberately revisable. Evidence that a peripheral system
can reproduce cue-eligibility decisions, intervention timing, sleep stability,
and behavioral outcomes relative to an EEG reference would weaken the rationale
for retaining EEG in later generations. Similarly, a replicated behavioral
advantage from event- or phase-aware control would justify increasing temporal
specificity, while failure of existing wearable EEG platforms to provide
adequate causal signal access or synchronization could justify dedicated
hardware development.

The first-generation pathway should therefore not be interpreted as a claim
that EEG is permanently superior to peripheral sensing, that stage-aware cueing
is physiologically optimal, or that existing hardware will necessarily satisfy
implementation requirements. It is a testable reference configuration chosen
because the current evidence provides the clearest traceable path from
established TMR physiology to an interpretable automated experiment. Detailed
algorithm design, numerical operating thresholds, device selection, deployment
format, and product validation remain subsequent research questions rather than
conclusions of the present synthesis.